\documentclass[12pt]{article}
\usepackage[utf8]{inputenc}
\usepackage[russian]{babel}
\usepackage{amsmath}
\usepackage{amssymb}
\usepackage{geometry}
\geometry{a4paper, margin=2cm}

\begin{document}

\section*{Задача 1.8}

\textbf{Условие:} Распознаватель $L=\{a^n b^n:n\geq 0\}$. Провести анализ сложности
в худшем и среднем случаях.

\section*{Программа}

\begin{verbatim}
alphabet = [ #, a, b]

tapes = [ right ]

START:
  # -> #R ON_LEFT_EDGE
  a -> aH ERROR
  b -> bH ERROR
//  #  ? ...  #
//     ^
ON_LEFT_EDGE:
  # -> #H STOP
  a -> #R GO_RIGHT
  b -> bH ERROR
//  # ... #
//     ^
GO_RIGHT:
  # -> #L ON_RIGHT_EDGE
  a -> aR GO_RIGHT
  b -> bR GO_RIGHT
//  # ... ? #
//        ^
ON_RIGHT_EDGE:
  # -> #H ERROR
  a -> aH ERROR
  b -> #L GO_LEFT
//  # ... #
//     ^
GO_LEFT:
  # -> #R ON_LEFT_EDGE
  a -> aL GO_LEFT
  b -> bL GO_LEFT
\end{verbatim}

\subsection*{Сложность}

Пусть $m$ - длина оставшегося слова при $ON\_LEFT\_EDGE$.
Если первый символ $a$, а последний - $b$
(при $m\geq 2$), выполняется итерация.

\begin{itemize}
  \item $ON\_LEFT\_EDGE$ - $1$ шаг
  \item $GO\_RIGHT$ - надо пройти $m-1$ оставшихся символов
   и перейти на правой границе $\#$ - $m$ шагов
  \item $ON\_RIGHT\_EDGE$ - 1 шаг
  \item $GO\_LEFT$ - надо пройти $m-2$ символа
  и перейти через левую границу - $m-1$ шага
\end{itemize}

Значит одна успешная итерация:
\[
  T(m)=1+m+1+(m-1)=2m+1
\]

Если активный участок непуст и итерация завершается ошибкой:
если первый символ - $b$ - ошибка за 1 шаг. Если первый символ
- $a$, машина стирает его, проходит вправо оставшиеся $m-1$ символов
и границу, после чего выполняет $ON\_RIGHT\_EDGE$.
Значит на неуспешную итерация - не более $m+2$ шагов.
Худшая ситуация получается, когда первый
символ - $a$, а последний символ - не $b$.

\subsection*{Сложность в худшем случае}

Пусть после успешных итераций (а их $(n-m)/2-1$ так как
удаляем по два символа на успешной итерации) 
остаётся $m$ символов.
Нам надо было пройти $START$, $(n-m)/2-1$ успешных итераций и завершить неуспешную итерацию.
Значит времени работы в худшем случае не больше:
\[
 T_n(m)=1+\sum_{j=0}^{(n-m)/2-1}\bigl(2(n-2j)+1\bigr)+\begin{cases}
 1,&m=0\\
 m+2,&m>0
 \end{cases}
\]
Для $m\geq3$:
\[
 T_n(m-2)-T_n(m)=(2m+1)+m-(m+2)=2m-1>0
\]
Для $m=2$:
\[
 T_n(0)-T_n(2)=5+1-4=2>0
\]
Значит худший случай будет при минимальном остатке:
$m=0$ для чётного $n$ и $m=1$ для нечётного $n$.

\paragraph{$n=2r$.}
Худшее слово: $a^r b^r$.
\[
 T_{\max}(2r)
 =1+\sum_{i=1}^{r}(4i+1)+1
 =2r^2+3r+2
\]

\paragraph{$n=2r+1$.}
Худшее слово: $a^{r+1}b^r$. После $r$ успешных итераций
остаётся $a$, на котором требуется $3$ перехода:
\[
 T_{\max}(2r+1)
 =1+\sum_{i=1}^{r}(4i+3)+3
 =2r^2+5r+4
\]

В итоге, во всех случаях:
\[
 T_{\max}(n)=\frac{n^2+3n+4}{2}
\]

\subsection*{Сложность в среднем случае}

Пусть $A_m$ - среднее число шагов от $ON\_LEFT\_EDGE$
до остановки для слова длины $m$. Тогда
\[
 A_0=1,\qquad A_1=\frac{1+3}{2}=2
\]
Для $m\geq2$:
\begin{enumerate}
  \item первый символ $b$ - ошибка за 1 шаг, вероятность $\frac12$
  \item первый символ $a$, последний символ $a$ - ошибка за $m+2$ шага, вероятность $\frac14$
  \item первый символ $a$, последний символ $b$ - успешная итерация за $2m+1$ шагов
   и переход к слову длины $m-2$ - $A_{m-2}$, вероятность $\frac14$
\end{enumerate}
Значит
\[
 A_m
 =\frac12\cdot1+\frac14(m+2)
   +\frac14(2m+1+A_{m-2})
 =\frac{3m+5+A_{m-2}}4
\]
Тогда
\[
A_0=1, A_1=2, A_2=3, A_3=4, A_4=5, \ldots
\]
Значит надо показать, что $A_m=m+1$. База это $A_0=1, A_1=2$.
По индукционному предположению $A_{m-2}=m-1$, значит
\[
 A_m=\frac{3m+5+(m-1)}4=m+1
\]
Тогда, $A_n=n+1$. Плюс начальный переход из $START$:
\[
 T_{\mathrm{avg}}(n)=1+A_n=n+2
\]

\end{document}